\section{Two dyadic compression mechanisms: integer cells and transport fibers}
\label{sec:two-compressions}
%======================================================================

Two exact identities in this program compress a M\"obius sum by M\"obius doubling.  Both are genuinely \emph{dyadic}: each pairs a source with its factor-$2$ multiple and uses the sign law $\mu(2m)=-\mu(m)$ on odd $m$.  They are nonetheless different theorems acting on different objects---one on consecutive integers inside a fixed cell, the other on the cofactor fibers of the transport term---and they must not be conflated.  Keeping them distinct is a permanent structural requirement, not a matter of exposition.

\subsection{Mechanism~I: four-slot cell doubling}

The first mechanism is the exact four-slot identity of \cref{thm:intro-dyadic-cell},
\[
 \bigl(\mu(4k+1),\mu(4k+2),\mu(4k+3),\mu(4k+4)\bigr)
 =\bigl(\mu(4k+1),-\mu(2k+1),\mu(4k+3),0\bigr),
\]
whose active content is the M\"obius \emph{doubling} law $\mu(4k+2)=-\mu(2k+1)$ together with $\mu(4k+4)=0$.  This is a factor-$2$ identity on \emph{consecutive integers} inside a fixed length-four cell.  It operates on the complete sum $M$ before any smooth/transport separation, and it is the compression referenced in the intake geometry and in the dyadic-cofactor cascade diagnostic of \cref{sec:numerics}.

\subsection{Mechanism~II: dyadic transport compression of parent/child packets}

The second mechanism lives entirely inside the transport channel and also uses the prime $2$, but now on the \emph{cofactor} of a channel rather than on consecutive integers.  Fix an odd cofactor $c$ and an odd upper prime $q$.  The parent channel $(c,q)$ and its doubled child $(2c,q)$ share the same transition index $q-1$ and have nested square-root entries $\lfloor\sqrt{cq}\rfloor\le\lfloor\sqrt{2cq}\rfloor$, while the doubling law
\begin{equation}
 \mu(2c\cdot q)=-\mu(c\cdot q)
 \qquad(c,q\ \text{odd})
 \label{eq:mobius-doubling-transport}
\end{equation}
gives them opposite source weights.  Writing $\mathcal P(c)$ for the finite transport packet of $(c,q)$, the parent splits into the boundary prefix between the two entry scales and the window it shares with the child,
\begin{equation}
 \mathcal P(c)=\mathcal B(c)+\mathcal C(c),
 \qquad
 \mathcal B(c)=\!\!\sum_{\lfloor\sqrt{cq}\rfloor\le t<\lfloor\sqrt{2cq}\rfloor}\!\!\mu(cq),
 \label{eq:packet-boundary-common}
\end{equation}
and the common window cancels identically against the child,
\begin{equation}
 \boxed{\ \mathcal P(c)+\mathcal P(2c)=\mathcal B(c).\ }
 \label{eq:packet-dyadic-cancel}
\end{equation}
Only the short boundary packet $\mathcal B(c)$, supported on the dyadic entry gap $\lfloor\sqrt{cq}\rfloor\le t<\lfloor\sqrt{2cq}\rfloor$, survives; every unmatched child is an explicit finite-horizon or born-smooth boundary, with no unclassified third case.

\subsection{The complete odd-annulus identity}

Applying the doubling involution \cref{eq:packet-dyadic-cancel} to every lower-cofactor fiber collapses each complete Möbius prefix onto its odd cofactors in a single dyadic boundary.  With $M(B)=\sum_{m\le B}\mu(m)$,
\begin{equation}
 \boxed{
 M(B)=\!\!\sum_{\substack{B/2<c\le B\\ c\ \mathrm{odd}}}\!\!\mu(c).
 }
 \label{eq:odd-annulus-identity}
\end{equation}
This is the complete odd dyadic annulus: no source is discarded, and at the manuscript endpoint it returns the full square-prefix value, $S_n=M(X_n)=\sum_{X_n/2<c\le X_n,\ c\ \mathrm{odd}}\mu(c)$.  Applied inside each prime fiber of the prime-first transport transfer \cref{eq:transport-primefirst}, with $B_q=\lfloor X_n/q\rfloor$, the entire transport term compresses to odd cofactors in the reciprocal dyadic boundary,
\begin{equation}
 \boxed{
 T_n=\!\!\sum_{\substack{R_n<q\le X_n\\ q\ \mathrm{prime}}}\
 \sum_{\substack{B_q/2<c\le B_q\\ c\ \mathrm{odd}}}\!\!\mu(c),
 }
 \label{eq:transport-dyadic}
\end{equation}
where every retained source $cq$ satisfies the exact canonical side conditions $c$ odd, $c<R_n<q$, $q=\Pplus(cq)$, and $cq\le X_n<2cq$.  The reindexing is finite Fubini followed by M\"obius doubling in every fiber; the paper transport coefficient $\mu(c)$ is the negative of the canonical source coefficient $\mu(cq)$.

\subsection{The permanent warning: high transport is only part of the annulus}

\begin{caution}[Transport is only part of the annulus]
\label{warn:transport-partial}
The complete odd dyadic annulus \cref{eq:odd-annulus-identity} equals the \emph{full} square-prefix Mertens value $S_n$; it contains every source through $X_n$, not only the transported ones.  The dyadic transport compression \cref{eq:transport-dyadic} and the prime-dilation reduction of \cref{sec:scale-transfer} act only on the \emph{transport} statistic $T_n$, i.e.\ on the one-large-prime population $\Pplus(m)>R_n$, which is a strict sub-part of that annulus.  They therefore do \emph{not} by themselves control the high-sector signal
\[
 S_n^{\mathrm{high}}(\Lambda)
 =\!\!\sum_{\substack{m\le X_n\\ |q_m^2-c_m^2|>2\Lambda\lfloor\sqrt m\rfloor}}\!\!\mu(m),
\]
whose height annulus $2\Lambda t<|q_m^2-c_m^2|\le2\Lambda(t+1)$ also contains born-smooth sources $q_m\le c_m$ belonging to $A_n$.  A bound on transport alone, on the complete annulus alone, or on the death process $\Delta D_t$ alone is insufficient: the smooth channel $A_n$ and the transport channel $T_n$ must be controlled \emph{jointly}, with their full signed interaction retained, because the exceptional aggregate cancellation observed numerically occurs precisely \emph{between} the born-smooth part of $A_n$ and $T_n$.  No transport-only, annulus-only, or shell-slice estimate closes the remaining inequality.
\end{caution}

\begin{remark}[Why the two mechanisms must not be merged]
\label{rem:no-merge}
Both mechanisms use the same prime and the same sign law $\mu(2m)=-\mu(m)$, so they cannot be told apart by their arithmetic alone; they differ only in what they act on.  Mechanism~I pairs \emph{consecutive integers} inside a fixed length-four cell and operates on the total sum $M$ before any smooth/transport split; Mechanism~II pairs a \emph{cofactor} with its double inside a single prime fiber of the transport term $T_n$, after that split.  Substituting one for the other silently changes the population and the object even though the prime is unchanged---exactly the confusion that the separate cell-mask and transport namespaces of the formalization are designed to forbid.  A four-slot cell scalar may never be used as a bound for a transport boundary packet $\mathcal B(c)$, or the converse.
\end{remark}

%======================================================================
